\section{Pollux Philosophy}~\label{sec:philosophy}

Any data language with way to specify a schema for a particular type of message
can have compatibility questions, simply put, does this encoded binary blob
conform to the known specification of this message type?

This project maintains a separation between the data description language and
the encoding format, enabling all parsing to be seen as two steps, one to parse
as much data from a binary blob as possible and another to validate it against
the data description. For formats like CBOR with CDDL, this is a natural fit
since CBOR is self-contained and any valid CBOR blob can be decoded to a value
and then compared against the CDDL specification. In that sense, the second step
for CBOR is really a validation step, while protobuf needs the information in
the descriptor to perform a series of type casting operations which can change
the values in the record before validation occurs.

However, this perspective isn't a natural fit for protobuf. Recall
Table~\ref{tab:tags}, and notice that encoding any of the variable width integer
fields in the proto language leads to the same tag value in the encoded
struct. While parsing one a \wtvar{} field, it is not possible to know
the final integer value being encoded without referencing some specification to
disambiguate between all the integer fields. However, that doesn't mean that
protobuf cannot be parsed without a descriptor to validate against. The
\texttt{protoscope} tool does this and is still able to provide some meaningful
information about an unknown protobuf blob. Some considerations about each tag
in the protoscope language are given below.

\begin{itemize}
\item All \texttt{varint} fields are unsigned integers which are then cast into
  the correct type according to the rules and functions in
  Section~\ref{sec:typ-rel}.
\item Tags (which would be types in \texttt{protoscope}, using the term to refer
  to the language rather then the tool. Use context to differentiate.) for
  \wtis{} and \wtil{} could only be interpreted as a list of bytes of
  the correct length. Casting to either a \texttt{fixed}, \texttt{sfixed} or
  floating point number would be part of the validation process.
\item \texttt{LEN} is the tricky case: its bytes could be a string, raw bytes,
  a packed list of numbers, or a nested message. \texttt{protoscope} guesses
  (\texttt{-all-fields-are-messages}) and can guess wrong. Pollux does not
  guess. The first step keeps every \texttt{LEN} payload as raw bytes, and the
  second step, which has the descriptor, decides what they are. For a nested
  message this means running both steps again on the payload, which is the loop
  in Figure~\ref{fig:parse-validate}.
\end{itemize}

Represented visually, the process for both CBOR+CDDL and Protobuf would look
like this:

\begin{figure}[H]
  \centering
  \begin{subfigure}[b]{\textwidth}
    \centering
    \begin{tikzpicture}
      \node[draw] (blob) at (0, 0) {Binary Blob};
      \node[draw] (ir) at (5, 0) {Low-Level Struct};
      \node[draw] (fin) at (10, 0) {Final Struct};

      \draw[-{LaTeX}] (blob) -- node[above]{Parsing} (ir);
      \draw[-{LaTeX}] (ir) -- node[above]{Validation} (fin);
    \end{tikzpicture}
    \caption{Progression from binary blob to fully parsed and validated struct.}
  \end{subfigure}

  \begin{subfigure}[b]{0.4\textwidth}
    \centering
    \begin{tikzpicture}
      \node[draw] (pp) at (0, 0) {P}; \node[draw] (pv) at (3, 0) {V};

      \draw[-{LaTeX}] (pp) edge[bend right] (pv); \draw[-{LaTeX}] (pv) edge[bend right] (pp);
    \end{tikzpicture}
    \caption{Protobuf parse validation loop for nested messages.}
  \end{subfigure}
  \hspace{1cm}
  \begin{subfigure}[b]{0.4\textwidth}
    \centering
    \begin{tikzpicture}
      \node[draw] (cp) at (0, 0) {P}; \node[draw] (cv) at (3, 0) {V};

      \draw[-{LaTeX}] (cp) -- (cv);
    \end{tikzpicture}
    \caption{CBOR parse validation flow. Tag 24 might require parsing again, but
    that is expected to be triggered by the host application at a later time.}
  \end{subfigure}

  \caption{Splitting the Protobuf and CDDL parse process into two steps, parsing
    and validating.}
  \label{fig:parse-validate}
\end{figure}

For Protobuf, the first step cuts the bytes into \emph{records}: a field number,
a wire type, and a raw payload. The wire type says where each payload ends, so
this needs no descriptor:
\[ \records : \mathrm{Bytes} \to \mathrm{Option}\ (\mathrm{List}\ (\id, w, \pay)). \]
The second step builds the reader's message one field at a time. The slot for
field $\id$ is $\assemble_{\dc[\id]}$ applied to $\restr{\records(\bs)}{\id}$, the
records with field number $\id$. A oneof group is handled as one unit, since a
record for one member clears the others. Records whose field number is unknown
to the reader, or whose wire type doesn't fit, are skipped. For Protobuf this
second step does more than validate: it converts numbers between types,
combines repeated records, and parses nested messages, and that is where all of
Protobuf's compatibility behaviour lives.

The compatibility relations never mention bytes. The semantic definitions of
Section~\ref{sec:framework} mention records, and only $\records$ and
$\encodes{\ty}$ mention bytes. Because each field's slot depends only on that
field's records, compatibility of whole messages follows from compatibility of
their fields.

Not every format has a first step. Avro needs the writer's schema to read
anything, and the format of Section~\ref{sec:inter-parse} cannot skip an unknown
field because its tags carry no wire type. Pollux's definition of compatibility
(Definition~\ref{def:compat}) does not depend on the split. The split is an
extra that makes field-by-field proofs possible when a format offers it.

% Local Variables:
% citar-bibliography: ("../pollux.bib")
% TeX-master: "../pollux.tex"
% TeX-command-default: "Make"
% End:
